Un électron ($m_e = \qty{9.109e-31}{kg}$, $e = \qty{1.602e-19}{C}$) entre avec une vitesse de \qty{\v{v}e6}{m/s} dans une région où règne un champ magnétique uniforme de \qty{\v{B}}{mT}. Sa vitesse fait un angle de \qty{\v{th}}{\degree} avec le champ, de sorte qu'il décrit une trajectoire en hélice.

\begin{enumerate}[label=\alph*)]
    \item Déterminez les composantes de la vitesse parallèle et perpendiculaire au champ.
    \item Déterminez le rayon de l'hélice.
    \item Déterminez le pas de l'hélice, c'est-à-dire la distance parcourue le long du champ pendant un tour.
\end{enumerate}
